% ==================================================
% 1. TIPOGRAFÍA Y FUENTES
% ==================================================
% Carga fuentes principales con detección automática:
% Garamond (preferente) → Times → TeX Gyre Termes
% Arial / Helvetica para sans-serif; Consolas / Courier para monoespaciado.

\IfFontExistsTF{EB Garamond}{%
  \setmainfont{EB Garamond}[
    UprightFont    = * Regular,
    BoldFont       = * Bold,
    ItalicFont     = * Italic,
    BoldItalicFont = * Bold Italic
  ]
}{%
  \IfFontExistsTF{Times New Roman}{\setmainfont{Times New Roman}}{%
    \setmainfont{TeX Gyre Termes}}}

\IfFontExistsTF{Arial}{\setsansfont{Arial}}{%
  \IfFontExistsTF{Helvetica}{\setsansfont{Helvetica}}{\setsansfont{TeX Gyre Heros}}}

\IfFontExistsTF{Consolas}{\setmonofont{Consolas}}{%
  \IfFontExistsTF{Courier New}{\setmonofont{Courier New}}{\setmonofont{TeX Gyre Cursor}}}

% ==================================================
% 2. IDIOMA Y TRADUCCIONES
% ==================================================
\usepackage[english,french,catalan,spanish]{babel}
% \selectlanguage{spanish}  % redundante: español es el último en la lista babel
\usepackage{csquotes}

\addto\captionsspanish{%
  \renewcommand{\contentsname}{Índice general}%
  \renewcommand{\listfigurename}{Índice de figuras}%
  \renewcommand{\listtablename}{Índice de tablas}%
  \renewcommand{\figurename}{Figura}%
  \renewcommand{\tablename}{Tabla}%
}

% ==================================================
% 3. DISEÑO DE PÁGINA Y NUMERACIÓN
% ==================================================
\usepackage{geometry}
\geometry{
  inner=3.5cm, outer=2.5cm,
  top=2.5cm, bottom=2.5cm
}

\setcounter{secnumdepth}{3} % hasta subsubsección
\setcounter{tocdepth}{2}    % hasta subsección en ToC

% ==================================================
% 4. TÍTULOS DE CAPÍTULOS Y SECCIONES
% ==================================================
\usepackage{titlesec}
\usepackage{setspace}

\makeatletter
% ----- Capítulos numerados -----
\def\@makechapterhead#1{%
 \vspace*{-20pt}
 {\parindent \z@ \centering \normalfont
 \Huge\bfseries \@chapapp\space \thechapter
 \par\nobreak\vskip 10pt
 {\singlespacing \LARGE \bfseries #1\par}\vskip 20pt}}

% ----- Capítulos sin numerar -----
\def\@makeschapterhead#1{%
 \vspace*{-20pt}
 {\parindent \z@ \centering \normalfont
 {\singlespacing \Huge \bfseries #1\par}\vskip 20pt}}

% ----- Redefinición general -----
\renewcommand{\@chapapp}{Capítulo}
\renewcommand{\thechapter}{\arabic{chapter}}
\def\@chapter[#1]#2{%
  \ifnum \c@secnumdepth >\m@ne
    \if@mainmatter
      \refstepcounter{chapter}%
      \addcontentsline{toc}{chapter}{%
        \protect\numberline{\thechapter.}#1}%
    \else
      \addcontentsline{toc}{chapter}{#1}%
    \fi
  \else
    \addcontentsline{toc}{chapter}{#1}%
  \fi
  \chaptermark{#1}%
  \addtocontents{lof}{\protect\addvspace{10\p@}}%
  \addtocontents{lot}{\protect\addvspace{10\p@}}%
  \if@twocolumn
    \@topnewpage[\@makechapterhead{#2}]%
  \else
    \ifx\relax#1\relax
      \@makeschapterhead{#2}%
    \else
      \@makechapterhead{#2}%
    \fi
    \@afterheading
  \fi}
\makeatother

% ----- Secciones y subniveles -----
\titleformat{\section}
 {\normalfont\Large\bfseries\singlespacing}{\thesection}{1em}{}
\titleformat{\subsection}
 {\normalfont\large\bfseries\singlespacing}{\thesubsection}{1em}{}
\titleformat{\subsubsection}
 {\normalfont\normalsize \bfseries\singlespacing}{\thesubsubsection}{1em}{}
\titleformat{\paragraph}[block]
 {\normalfont\normalsize\itshape\singlespacing}{}{0pt}{}
% Espaciado antes/después de secciones y subniveles (reducción ~40% sobre
% los valores por defecto de titlesec). No afecta al título de capítulo,
% que se gobierna desde \@makechapterhead más arriba en este preámbulo.
\titlespacing*{\section}{0pt}{2ex plus 0.5ex minus 0.2ex}{1ex plus 0.2ex}
\titlespacing*{\subsection}{0pt}{1.6ex plus 0.4ex minus 0.2ex}{0.8ex plus 0.2ex}
\titlespacing*{\subsubsection}{0pt}{1.4ex plus 0.3ex minus 0.2ex}{0.6ex plus 0.2ex}
\titlespacing*{\paragraph}{0pt}{1ex plus .2ex minus .2ex}{0.5ex}

% ==================================================
% 5. ENCABEZADOS Y PIES DE PÁGINA
% ==================================================
\usepackage{fancyhdr}
\pagestyle{fancy}
\fancyhf{}
\fancyfoot[LE,RO]{\thepage}
\renewcommand{\headrulewidth}{0pt}
% Páginas de apertura de capítulo (estilo 'plain'): número en margen exterior
\fancypagestyle{plain}{%
  \fancyhf{}%
  \fancyfoot[LE,RO]{\thepage}%
  \renewcommand{\headrulewidth}{0pt}%
}
\usepackage{emptypage}

% ==================================================
% 6. ÍNDICES Y LISTAS (ToC, Figuras, Tablas)
% ==================================================
\usepackage{tocloft}
\renewcommand{\cfttoctitlefont}{\hspace*{\fill}\Huge\bfseries}
\renewcommand{\cftloftitlefont}{\hspace*{\fill}\Huge\bfseries}
\renewcommand{\cftlottitlefont}{\hspace*{\fill}\Huge\bfseries}
\renewcommand{\cftaftertoctitle}{\hspace*{\fill}}
\renewcommand{\cftafterloftitle}{\hspace*{\fill}}
\renewcommand{\cftafterlottitle}{\hspace*{\fill}}

\renewcommand{\cftchapfont}{\bfseries}
\renewcommand{\cftchappagefont}{\bfseries}
\renewcommand{\cftchapnumwidth}{2em}
\renewcommand{\cftchapindent}{0pt}
\renewcommand{\cftsecfont}{\normalsize}
\renewcommand{\cftsubsecfont}{\normalsize}
\setlength{\cftbeforesecskip}{4pt}

% ==================================================
% 7. FIGURAS, TABLAS Y CAPTIONS
% ==================================================
\usepackage{caption}
% labelfont: normalsize (no large) para evitar etiquetas visualmente pesadas
\captionsetup{labelfont={normalsize,bf},textfont=normalfont,labelsep=colon}

\usepackage{rotating}
\usepackage{eso-pic} % marca de agua centrada en página (declaración jurada)
\usepackage{threeparttable,placeins,float,booktabs,longtable,tabularx,multirow,array,ragged2e}
\usepackage[table]{xcolor}
% Asegura que las notas de tabla (\tnote) usen \normalsize.
\AtBeginEnvironment{threeparttable}{\normalsize}

% ==================================================
% 8. MACROS EDITORIALES (etoolbox ANTES de microtype)
% ==================================================
\usepackage{etoolbox}

% ==================================================
% 9. TIPOGRAFÍA Y JUSTIFICACIÓN
% ==================================================
\usepackage{microtype}
\tolerance=1000
\emergencystretch=3em
% \sloppy % <- ELIMINADO: en conflicto con ajustes manuales

% Interlineado 1.3 (aprox. 1.5 en Word)
\setstretch{1.3}
% Espaciado entre párrafos uniforme: parskip rígido (sin elasticidad) y
% \raggedbottom para que LaTeX deje páginas ligeramente más cortas antes
% que estirar el espacio inter-párrafo, evitando el efecto de huecos
% desiguales entre párrafos en páginas con flotantes o tablas.
\setlength{\parindent}{0pt}
\setlength{\parskip}{0.6\baselineskip}
\raggedbottom
% Notas al pie a espacio simple
\AtBeginEnvironment{footnote}{\singlespacing}
% Tablas a espacio simple dentro de las celdas (convención académica): el
% \setstretch{1.3} del cuerpo no se hereda en el contenido tabular, por lo
% que las celdas con texto envuelto en varias líneas mantienen una
% separación interna compacta. Cubre los entornos que kableExtra puede
% generar según latex_options (tabular base, longtable para repeat_header,
% tabu/longtabu para striped, tabularx para ancho total).
% Se usa \setstretch{1} en lugar de \singlespacing porque éste último
% incluye \vskip\baselineskip al final, lo que es ilegal dentro del
% preámbulo de tabular (sobre todo cuando va envuelto en \resizebox por
% scale_down de kableExtra).
\AtBeginEnvironment{tabular}{\setstretch{1}}
\AtBeginEnvironment{tabularx}{\setstretch{1}}
\AtBeginEnvironment{longtable}{\setstretch{1}}
\AtBeginEnvironment{tabu}{\setstretch{1}}
\AtBeginEnvironment{longtabu}{\setstretch{1}}
% Altura de fila: 20% por encima del default. Separa visualmente las filas
% en tablas con celdas multilínea (donde el interlineado compacto interno
% antes confundía las fronteras entre categorías consecutivas), sin afectar
% el interlineado dentro de cada celda, que sigue siendo \setstretch{1}.
\renewcommand{\arraystretch}{1.2}

% ==================================================
% 10. HIPERVÍNCULOS Y METADATOS PDF
% ==================================================
\usepackage{hyperref}
\usepackage{bookmark}
\hypersetup{
  bookmarksnumbered = true,
  colorlinks        = true,
  linkcolor         = black,
  citecolor         = black,
  urlcolor          = darkgray,
  breaklinks        = true,
  pdfpagelabels     = true,
  pdftitle          = {La vulnerabilidad social como fenómeno estructural y contextual},
  pdfauthor         = {Juan Antonio Romero Crespo},
  pdfsubject        = {Tesis doctoral},
  pdfkeywords       = {vulnerabilidad, riesgo, IVEC, ciencias sociales}
}

% ==================================================
% 11. MACROS — Glosario y abreviaturas
% ==================================================
% Conserva jerarquías, añade espaciado uniforme y mantiene etiquetas para hipervínculos internos.

\newcommand{\referenceletter}[2]{%
 \vspace{0.1em}%
 \par\noindent{\Large\bfseries #1}\par
 \vspace{0.1em}%
 \noindent\rule{\textwidth}{0.4pt}\par
 \vspace{0.3em}%
 \pdfbookmark[1]{#1}{gloss-#2}%
}

\newcommand{\tesisglossary}[2]{%
 \vspace{0.3em}%
 \par\noindent{\textbf{#1}}\par
 \vspace{0.2em}%
 \phantomsection%
 \label{gloss-#2}%
 \pdfbookmark[2]{#1}{gloss-#2}%
}

% ==================================================
% 12. BIBLIOGRAFÍA — EVITAR PARTICIÓN DE ENTRADAS
% ==================================================
% Pandoc-citeproc genera la bibliografía dentro del entorno CSLReferences.
% Cada entrada es un párrafo independiente; \interlinepenalty=10000 impide
% que LaTeX parta una entrada entre el final de una página y el inicio
% de la siguiente. El valor 10000 es el máximo que puede asignarse sin
% forzar un overfull \vbox.
% etoolbox ya está cargado en la sección 8 de este preámbulo.
\AtBeginEnvironment{CSLReferences}{\interlinepenalty=10000}

% ==================================================
% 13. APÉNDICES COMO CAPÍTULOS (bookdown: encabezados con #)
% ==================================================
\makeatletter
\newcommand{\startappendiceschapter}{
  \cleardoublepage
  \appendix

  % Etiqueta del capítulo (Capítulo -> Apéndice)
  \renewcommand{\@chapapp}{Apéndice}

  % Numeración alfabética de capítulos
  \renewcommand{\thechapter}{\Alph{chapter}}
  \renewcommand{\theHchapter}{\thechapter}

  % Figuras y tablas por apéndice (A.1, A.2...)
  \renewcommand{\thefigure}{\thechapter.\arabic{figure}}
  \renewcommand{\thetable}{\thechapter.\arabic{table}}

  % Reinicios (opcional, pero recomendable)
  \setcounter{chapter}{0}
  \setcounter{figure}{0}
  \setcounter{table}{0}
}
\makeatother
